\section{Day 19: (Nov.\ 24, 2026)}
Let $L/K$ be number fields, let $\kp \subset \SO_K$ be prime, and pick $\kq_1, \dots, \kq_r \subset \SO_L$ primes over $\kp$, where we denote $k_{\kq_i} = \SO_L / \kq_i$. For $x \in \SO_L$,
\[ \ol{\Tr_{L/K}(x)} = \sum_{i=1}^r \Tr_{k_{\kq_i}/k_\kp}(\ol x_i). \]
Note that $L/K$ lies in $\SO_K/\kp = k_\kp$, and we write $\ol x_i = x \mod \kq_i$ in $k_{\kq_i}$.

\begin{theorem} \label{thm:rational_prime_discriminant_ramification}
    Let $L$ be a number field, and $p \in \ZZ$ be a rational prime. If $p \mid \opname{disc}(L)$, then $p$ is ramified in $L$.
\end{theorem}
\begin{proof}
    For the first step, we reduce to the case where $L/\QQ$ is Galois, and consider $M/\QQ$, the Galois closure of $L$. Recall that $M$ is the composite field of all embeddings $L \hookrightarrow \CC$. We saw that if $p$ is unramified in $L$, then it's also unramified in $M$, and moreover,
    \[ \opname{disc}(L) \mid \opname{disc}(K). \]
    It is enough to prove that, if $p \mid \opname{disc}(M)$, then $p$ is ramified in $M$. Thus, without loss of generality, suppose $L/\QQ$ is Galois.

    For the second step, suppose $p \mid \opname{disc}(L)$, and suppose $p$ is unramified. We have that
    \[ (p) = p \SO_L = \kq_1 \dots \kq_r \]
    with $\kq_1, \dots, \kq_r$ distinct primes. Since $p \mid \opname{disc}(L)$, there exists $a \in \SO_L \setminus p \SO_L$ such that $\Tr(ab) \in p \ZZ$ for all $b \in \SO_K$. {\color{amethyst} This is because the discriminant is the determinant of bilinear forms on $\ZZ^n$ ($n = \abs{L : \QQ}$), given by $(a, b) = \Tr_{L/\QQ}(a \cdot b)$.} $p \mid \opname{disc}(L)$ implies that this bilinear form is singular mod $p$, so there exists $a$ such that $\ol a \in (\ZZ/p)^a$ is nonzero, and $(\ol a, \ol b) = 0$ in $\ZZ/p$, i.e., there exists $a \in \SO_L \setminus p \SO_L$ such that $p \mid \Tr(ab)$ for all $b \in \SO_K$. We have $p \SO_L = \kq_1 \dots \kq_r$ with $a \notin p \SO_L$, so there exists $\kq_i$ such that $a \notin \kq_i$. Let $D = D(\kq_i \mid p)$. $p$ is unramified because its inertia group is trivial (i.e., $\{e\}$), so
    \[ D = \Gal(k_{p_i} / \FF_p). \]
    We claim that the trace pairing for $k_{\kq_i} / \FF_p$ is non-degenerate. There exists $\ol b \in k_{\kq_i}$ such that $\Tr(\ol a \cdot \ol b) \neq 0$ {\color{amethyst}($\smash{\ol a \cdot \ol b \in \FF_p}$)}, for which the Chinese remainder theorem yields the existence of some $b \in \SO_L$ such that $\smash{b \equiv \ol b}$ modulo $\kq_i$, and $b \equiv 0$ modulo $\kq_j$ for $j \neq i$. Applying the trace identity, we get
    \[ 0 = \ol{\Tr_{K/\QQ}(ab)} = \sum_{\ell=1}^r \Tr_{k_{\kq_\ell}/\FF_p}\bigl((\ol{ab})_\ell\bigr) = \Tr_{k_{\kq_i} / \FF_p}(\ol a \cdot \ol b) \neq 0, \]
    by our choice of $\ol b_i$, which is a contradiction.

    We now talk about non-degeneracy of the trace pairing for finite fields. We claim that $K/\FF_\kq$ \\ has a non-degenerate trace pairing. Let $a \in k \setminus \{0\}$. We need to show that there exists $b \in k$ such that
    \[ \Tr_{k/\FF_p}(a \cdot b) \neq 0. \]
    We have that $\Gal(k/\FF_q) = C_n$, where $n$ is the degree, generated by $x \mapsto x^q$. We also have
    \[ \Tr(x) = \sum_{\sigma \in \Gal(k/\FF_q)} \sigma(x) = x + x^q + \dots + x^{q^{n-1}}, \]
    which has degree $q^{n-1} < \abs{k} = q^n$. Thus, there exists $e \in k$ such that $\tr9e) = f(e) \neq 0$, and we may now take $b = a\inv c$. Then $\Tr(ab) = \Tr(e) \neq 0$.
\end{proof}

\newpage
We now talk about finiteness of class groups. Our goal is to show the following theorem.
\begin{theorem*} % \label{thm:class_group_number_field}
    Let $K$ be a number field. The class groups
    \[ \opname{Cl}(K) = \{I \subset \SO_K \text{ nonzero}\} / \sim, \]
    where $I \sim J$ if there exists $\alpha, \beta \in \SO_K$ such that $\alpha I = \beta J$. We saw that $\opname{Cl}(K)$ is a group.
\end{theorem*}
We will not do the proof now, but we will do it on Thursday. Here is a brief outline. We saw, that if $\opname{Cl}(K) = 0$, then $\SO_K$ is a PID.

We want to see, to what extent, the field norm (fails to) give the structure of a Euclidean domain. Try to use $\abs{N_{K/\QQ}(x)}$ as the function.
\begin{lemma*}
    Let $K$ be a number field. There exists an integer $M \geq 1$ such that, for all $\alpha, \beta \in \SO_K$, there exists $w, r \in \SO_K$ and integer $t \in \{1,\dots, M\}$ such that $t \alpha = w \cdot \beta + r$ and $\abs{N_{K/\QQ}(r)} \leq \abs{N_{K/\QQ}(p)}$.
\end{lemma*}

With this, we may prove the theorem by taking $M \geq 1$ as in the lemma, and see that $\SO_K/\left<M!\right>$ is finite. The set $S = \{I \subset \SO_K \mid M! \in I\}$ is finite.